\section{Submission Scope and Novelty Statement}

This manuscript addresses the ING-INF/05 area through a systems-oriented contribution at the intersection of distributed software systems, embedded computing, and secure cyber-physical infrastructures. The paper does not claim novelty in WebAssembly, Kubernetes, TLS, or CBOR as isolated technologies. Its contribution lies instead in the principled integration of these elements into a coherent and reproducible platform that closes an implementation gap frequently acknowledged in the literature but less often documented end to end: how declarative cloud-side intent can be made operational on constrained embedded devices without sacrificing secure attachment, control-plane visibility, and state fidelity.

The novelty of RETROSPECT is threefold. First, it proposes a Kubernetes-native resource model in which edge devices, gateways, and deployable applications are represented as first-class CRDs rather than ad hoc metadata external to the cluster. Second, it introduces a gateway-centric southbound control strategy that makes it feasible to retain Kubernetes-like orchestration semantics even when target endpoints cannot directly host container runtimes or Kubernetes agents. Third, it documents and resolves a non-trivial consistency problem between observed transport state and reported device status, showing that reliable edge supervision requires careful reconciliation policies rather than only a correct enrollment protocol.

Compared with isolated runtime demonstrations, RETROSPECT contributes a complete control-plane and data-plane integration, including controller logic, gateway-mediated secure enrollment, CRD status convergence, application life-cycle signaling, deployment reproducibility, and experimentally validated operational behavior. This positions the work as a systems-engineering contribution suited to an IEEE Transactions venue where methodological rigor, implementation completeness, and reproducibility are valued alongside architectural novelty.
